% SPDX-FileCopyrightText: © 2026 Idiap Research Institute <contact@idiap.ch>
%
% SPDX-FileContributor: Andre Freitas <andre.freitas@idiap.ch>
% Neuro-symbolic AI Group
%
% SPDX-FileContributor: Danilo Gusicuma <danilo.gusicuma@idiap.ch>
% Neuro-symbolic AI Group
%
% SPDX-License-Identifier: MIT
\documentclass[11pt]{article}

\usepackage[margin=1in]{geometry}
\usepackage[T1]{fontenc}
\usepackage[utf8]{inputenc}
\usepackage{lmodern}
\usepackage{microtype}
\usepackage{amsmath,amssymb,amsthm,mathtools}
\usepackage{enumitem}
\usepackage{booktabs}
\usepackage{array}
\usepackage{hyperref}
\usepackage{xcolor}

\hypersetup{
  colorlinks=true,
  linkcolor=blue!50!black,
  citecolor=blue!50!black,
  urlcolor=blue!50!black
}

\newtheorem{theorem}{Theorem}[section]
\newtheorem{lemma}[theorem]{Lemma}
\newtheorem{proposition}[theorem]{Proposition}
\newtheorem{corollary}[theorem]{Corollary}
\newtheorem{definition}[theorem]{Definition}
\newtheorem{remark}[theorem]{Remark}
\newtheorem{example}[theorem]{Example}

\newcommand{\one}{\mathbf 1}
\newcommand{\C}{\mathbb C}
\newcommand{\R}{\mathbb R}
\newcommand{\eps}{\varepsilon}
\newcommand{\ip}[2]{\langle #1,#2\rangle}
\newcommand{\diag}{\operatorname{diag}}
\newcommand{\Ran}{\operatorname{Ran}}
\newcommand{\Ker}{\operatorname{Ker}}
\newcommand{\OS}{\mathrm{OS}}
\newcommand{\YM}{\mathrm{YM}}
\newcommand{\Had}{\circ}
\newcommand{\odiv}{\oslash}
\newcommand{\snull}{S_{\mathrm{null}}}
\newcommand{\norm}[1]{\lVert #1\rVert}
\newcommand{\Span}{\operatorname{span}}
\newcommand{\Herm}{\operatorname{Herm}}
\newcommand{\Tr}{\operatorname{Tr}}
\newcommand{\rank}{\operatorname{rank}}
\newcommand{\Ree}{\operatorname{Re}}

\title{A Finite Schur--Feshbach Consistency Test for Normalized Reflected Wilson Gram Blocks}
\author{}
\date{}

\begin{document}
\maketitle

\begin{abstract}
Reflection positivity is tested on finite Gram matrices before it becomes an infinite-dimensional reconstruction problem.  This paper proves a finite Schur--Feshbach criterion for Hermitian matrix families whose leading term is a rank-one reflected-vacuum block.  If
\[
q_\eps=\one\one^*+\eps H+\eps^2K+o(\eps^2),
\]
then the correct quotient-space expansion after eliminating the rank-one direction is
\[
F_\eps
=\eps PHP+
\eps^2P\bigl(K-n^{-1}Hee^*H\bigr)P+o(\eps^2),
\qquad e=n^{-1/2}\one,
\qquad P=I-ee^*.
\]
Thus the second-order test is not the naive compression \(PKP\), but the Schur--Feshbach corrected form \(P(K-n^{-1}Hee^*H)P\).  The result gives necessary finite positivity conditions, strict sufficient margins, and a recursive null-sector test.  The intended use is as a finite consistency check inside Wilson-loop/Osterwalder--Schrader constructions, where normalized reflected Wilson blocks may approach a rank-one collision limit.  The claim is deliberately finite: the paper does not construct Yang--Mills theory and does not prove a mass gap.  It identifies the finite coefficient data that any same-theory Wilson/OS realization must supply and respect before such blocks can be used in a reconstruction argument.
\end{abstract}

\tableofcontents

\section{Introduction and claim boundary}

The Osterwalder--Schrader approach to quantum field theory constructs a real-time Hilbert-space theory from Euclidean correlation functions satisfying, among other conditions, reflection positivity \cite{OsterwalderSchrader1973,OsterwalderSchrader1975}.  In gauge theory, Wilson loops are natural gauge-invariant observables, and finite families of Wilson observables give finite reflected Gram matrices.  A Euclidean reconstruction programme must keep such finite matrices positive while removing regulators and passing to larger observable classes.

This paper isolates one finite step in that programme.  It concerns normalized reflected Wilson Gram blocks that approach a rank-one limit.  The finite situation is simple to state: a Hermitian matrix family has the expansion
\begin{equation}
  q_\eps=\one\one^*+\eps H+\eps^2K+o(\eps^2).
  \label{eq:intro-expansion}
\end{equation}
The leading matrix \(\one\one^*\) has one positive direction and kernel \(\one^\perp\).  Reflection-positivity information therefore lives in a singular quotient limit.  The paper proves that the correct second-order quotient coefficient is not \(PKP\).  One must first remove the rank-one direction by a Schur complement, producing the correction
\begin{equation}
  -\,n^{-1}PHee^*HP,
  \qquad e=n^{-1/2}\one,
  \qquad P=I-ee^*.
  \label{eq:intro-correction}
\end{equation}

The result has three levels, which should not be conflated.
\begin{enumerate}[label=(L\arabic*)]
  \item \textbf{Finite theorem.}  For any finite Hermitian family satisfying \eqref{eq:intro-expansion}, positivity near \(\eps=0\) forces \(PHP\succeq0\) on \(\one^\perp\) and, on \(\Ker(PHP)\), forces the corrected form \(P(K-n^{-1}Hee^*H)P\succeq0\).  Strict versions give sufficient local positivity margins.
  \item \textbf{Wilson-block application.}  To apply the theorem to QFT, the matrices \(q_\eps\) must be produced by a specified regulated Wilson observable algebra, a specified reflection map, and a same-theory normalization.  Entrywise quotienting is not automatically positivity-preserving.
  \item \textbf{Continuum reconstruction.}  Passing from finite regulated blocks to a continuum Yang--Mills theory with a mass gap requires additional uniform estimates, density/reconstruction arguments, and spectral information.  Those are outside the theorem proved here.
\end{enumerate}

The contribution is therefore not a negative claim about Yang--Mills.  It is a positive finite tool: a reader-checkable Schur--Feshbach test that tells a Wilson/OS construction exactly which coefficient forms must be computed, what sign conditions matter, and what transfer obligations remain.

\section{Finite reflection positivity and the quotient-block problem}

Let \(\Theta\) denote Euclidean time reflection, and let \(\mathcal A_+\) be a positive-time algebra of observables.  For a Euclidean Schwinger functional \(S\), reflection positivity says that, for every finite family \(F_1,\ldots,F_n\in\mathcal A_+\),
\begin{equation}
  \sum_{i,j=1}^n \overline{c_i}c_j\,S\bigl((\Theta F_i)^*F_j\bigr)\ge 0,
  \qquad c\in\C^n.
  \label{eq:os-pos}
\end{equation}
Equivalently, the reflected Gram matrix
\begin{equation}
  N_{ij}=S\bigl((\Theta F_i)^*F_j\bigr)
  \label{eq:reflected-gram}
\end{equation}
is positive semidefinite.

For Wilson observables one often has regulator-dependent one-body, perimeter, cusp, collar, or boundary factors.  A normalized finite block may be written schematically as
\begin{equation}
  q_\eps=\widehat N_\eps \odiv \widehat B_\eps,
  \label{eq:quotient}
\end{equation}
where \(\odiv\) denotes entrywise division.  This notation is only schematic until the normalization is realized in the same regulated theory.  If \(\widehat B_\eps\) is an individual one-body factor, the normalization may be an invertible diagonal congruence and therefore preserves positivity.  If \(\widehat B_\eps\) contains pair-dependent boundary data, entrywise division need not preserve positivity.  A Wilson/OS application must therefore prove that \(q_\eps\) is itself an allowed same-theory reflected block, or prove an equivalent positivity-preserving representation.

The finite problem studied here begins after that input is supplied.  Assume that the normalized block is Hermitian and has a rank-one collision expansion
\begin{equation}
  q_\eps\longrightarrow \one\one^*.
  \label{eq:rank-one-limit}
\end{equation}
The common vector \(\one\) is the leading reflected-vacuum direction of the finite block.  The nontrivial test is on \(\one^\perp\), but the rank-one direction cannot simply be ignored: it couples to \(\one^\perp\) at first order and contributes back at second order.

\section{The finite Schur--Feshbach theorem}

\subsection{Set-up}

Let \(q_\eps\in\Herm_n(\C)\) satisfy
\begin{equation}
  q_\eps=\one\one^*+\eps H+\eps^2K+R_\eps,
  \qquad \norm{R_\eps}=o(\eps^2),
  \label{eq:main-expansion}
\end{equation}
with \(H=H^*\), \(K=K^*\).  Set
\begin{equation}
  e=n^{-1/2}\one,
  \qquad P=I-ee^*.
  \label{eq:e-P}
\end{equation}
Then \(P\) is the orthogonal projection onto \(\one^\perp\), and \(\one\one^*=nee^*\).

\subsection{Main finite theorem}

\begin{theorem}[Finite Schur--Feshbach tangent-cone criterion]
\label{thm:finite-certificate}
Let \(q_\eps\) satisfy \eqref{eq:main-expansion}.  Define
\begin{equation}
  \lambda_\eps=e^*q_\eps e,
  \qquad
  F_\eps=Pq_\eps P-Pq_\eps e\,\lambda_\eps^{-1}\,e^*q_\eps P.
  \label{eq:F-eps}
\end{equation}
Then \(\lambda_\eps=n+O(\eps)>0\) for sufficiently small \(\eps>0\), and
\begin{equation}
  q_\eps\succeq0
  \quad\Longleftrightarrow\quad
  F_\eps\succeq0.
  \label{eq:schur-equivalence}
\end{equation}
Moreover,
\begin{equation}
  F_\eps=
  \eps A+\eps^2C+o(\eps^2),
  \qquad
  A=PHP,
  \qquad
  C=P\bigl(K-n^{-1}Hee^*H\bigr)P.
  \label{eq:A-C}
\end{equation}
Consequently, if \(q_\eps\succeq0\) for all sufficiently small \(\eps>0\), then
\begin{equation}
  A\succeq0\quad\text{on }\one^\perp.
  \label{eq:first-condition}
\end{equation}
If \(\Pi\) is the orthogonal projection in \(\one^\perp\) onto \(\Ker A\), then also
\begin{equation}
  \Pi C\Pi\succeq0\quad\text{on }\Ker A.
  \label{eq:second-condition}
\end{equation}
Equivalently, either of the following finite conditions rules out positivity for all sufficiently small \(\eps>0\):
\begin{align}
  &\exists v\in\one^\perp:
    \quad v^*PHPv<0,
    \label{eq:first-negative-direction}\\
  &\exists v\in\Ker(PHP):
    \quad v^*P\bigl(K-n^{-1}Hee^*H\bigr)Pv<0.
    \label{eq:second-negative-direction}
\end{align}
\end{theorem}

\begin{proof}
Because \(q_0=nee^*\),
\[
  \lambda_\eps=e^*q_\eps e=n+\eps e^*He+O(\eps^2),
\]
so \(\lambda_\eps>0\) for small \(\eps\).  With respect to the splitting \(\C^n=\C e\oplus P\C^n\), the upper-left block of \(q_\eps\) is the positive scalar \(\lambda_\eps\).  The scalar Schur complement criterion gives \eqref{eq:schur-equivalence} with \(F_\eps\) as in \eqref{eq:F-eps}.

Next,
\[
  Pq_\eps P=\eps PHP+\eps^2PKP+o(\eps^2),
\]
because \(P(\one\one^*)P=0\).  Also,
\[
  Pq_\eps e=\eps PHe+\eps^2PKe+o(\eps^2),
  \qquad
  \lambda_\eps^{-1}=n^{-1}+O(\eps).
\]
Hence
\begin{align*}
  Pq_\eps e\,\lambda_\eps^{-1}\,e^*q_\eps P
  &=\bigl(\eps PHe+O(\eps^2)\bigr)
    \bigl(n^{-1}+O(\eps)\bigr)
    \bigl(\eps e^*HP+O(\eps^2)\bigr)\\
  &=\eps^2 n^{-1}PHee^*HP+o(\eps^2).
\end{align*}
Substitution gives \eqref{eq:A-C}.

If \(q_\eps\succeq0\), then \(F_\eps\succeq0\).  For \(v\in\one^\perp\),
\[
  0\le v^*F_\eps v
  =\eps v^*Av+\eps^2v^*Cv+o(\eps^2).
\]
Dividing by \(\eps\) and taking \(\eps\downarrow0\) gives \(v^*Av\ge0\).  If \(v\in\Ker A\), then the order-\(\eps\) term vanishes; dividing by \(\eps^2\) gives \(v^*Cv\ge0\).  This proves the two necessary conditions, and the negative-direction statements are their contrapositives.
\end{proof}

\begin{proposition}[Strict second-order validation]
\label{prop:sufficiency}
Assume \(A\succeq0\) on \(\one^\perp\).  Let \(\Pi\) be the projection onto \(\Ker A\), and set \(Q=P-\Pi\).  Suppose there are constants \(\alpha,\beta>0\) such that
\begin{equation}
  QAQ\succeq \alpha Q,
  \qquad
  \Pi C\Pi\succeq \beta \Pi.
  \label{eq:strict-margins}
\end{equation}
Then \(F_\eps\succeq0\), hence \(q_\eps\succeq0\), for all sufficiently small \(\eps>0\).
\end{proposition}

\begin{proof}
On \(Q\C^n\),
\[
  QF_\eps Q=\eps QAQ+O(\eps^2)\succeq \tfrac12\eps\alpha Q
\]
for small \(\eps\), so this block is positive and invertible.  Since \(A\Pi=0=\Pi A\),
\[
  \Pi F_\eps Q=\eps^2\Pi C Q+o(\eps^2).
\]
The Schur complement of the positive \(Q\)-block is
\begin{align*}
  \Pi F_\eps\Pi
  -\Pi F_\eps Q\,(QF_\eps Q)^{-1}\,QF_\eps\Pi
  &=\eps^2\Pi C\Pi+o(\eps^2)-O(\eps^3)\\
  &\succeq \tfrac12\eps^2\beta\Pi
\end{align*}
for small \(\eps\).  Therefore \(F_\eps\succeq0\), and Theorem \ref{thm:finite-certificate} gives \(q_\eps\succeq0\).
\end{proof}

If \(\Pi C\Pi\) is semidefinite rather than strictly positive, the same Schur--Feshbach reduction can be repeated on its null sector.  In finite dimension, this gives a recursive tangent-cone algorithm: the first negative reduced coefficient is a genuine finite loss-of-positivity certificate, while strict positivity at a reduced stage gives a local margin.

\section{A minimal example showing why \texorpdfstring{\(PKP\)}{PKP} is insufficient}
\label{sec:minimal-example}

The following two-dimensional example is not a gauge-theory calculation.  Its purpose is to show, in the smallest possible setting, that the Schur--Feshbach correction can change the sign of the second-order test.

\begin{example}[Positive \(PKP\), negative corrected form]
\label{ex:naive-PKP-fails}
Let \(n=2\), let
\[
  e=2^{-1/2}(1,1)^T,
  \qquad
  f=2^{-1/2}(1,-1)^T,
  \qquad
  P=ff^*.
\]
Define
\begin{equation}
  H=ef^*+fe^*,
  \qquad
  K=\frac14 ff^*.
  \label{eq:toy-HK}
\end{equation}
In the orthonormal basis \((e,f)\), the matrix family
\[
  q_\eps=\one\one^*+\eps H+\eps^2K
\]
has the block form
\begin{equation}
  q_\eps
  =
  \begin{pmatrix}
    2 & \eps\\
    \eps & \eps^2/4
  \end{pmatrix}_{(e,f)}.
  \label{eq:toy-block}
\end{equation}
The naive compressed second-order coefficient is positive:
\[
  f^*PKPf=\frac14.
\]
The corrected coefficient is
\begin{align*}
  f^*P\bigl(K-2^{-1}Hee^*H\bigr)Pf
  &=\frac14-\frac12|f^*He|^2
  =\frac14-\frac12
  =-\frac14.
\end{align*}
Indeed,
\[
  \det q_\eps=2\cdot\frac{\eps^2}{4}-\eps^2=-\frac{\eps^2}{2}<0
\]
for every \(\eps>0\).  Thus \(q_\eps\) is indefinite even though the naive \(PKP\) test is positive.  The missing term is exactly the coupling through the rank-one direction \(e\).
\end{example}

This example is the finite warning that motivates the Wilson/OS use of the theorem.  A normalized Wilson quotient calculation that reports only \(PKP\) has not yet tested the correct second-order tangent cone.

\section{Same-theory transfer certificate}
\label{sec:same-theory}

The finite theorem becomes a QFT statement only after a regulated theory supplies its input matrices.  This section states the transfer certificate required for a Wilson/OS application.

Let \(\Lambda_a\) be a finite Euclidean lattice with spacing \(a\), compact gauge group \(G\), lattice gauge measure \(d\mu_{a}\), and reflection \(\Theta\) across a time-zero hyperplane.  Let \(\mathcal A_{+,a}\) be the positive-time gauge-invariant Wilson algebra.  A Wilson observable associated with a loop or path system \(\gamma\) in representation \(R\) has the form
\begin{equation}
  W_{\gamma,R}(U)=\frac{1}{d_R}\Tr_R\prod_{\ell\in\gamma}^{\rightarrow} U_\ell,
  \label{eq:wilson-observable}
\end{equation}
with the usual path ordering.  For a finite family \(F_i^{a,\eps}\in\mathcal A_{+,a}\), the reflected Gram block is
\begin{equation}
  N_{a,\eps}(i,j)
  =\int (\Theta F_i^{a,\eps})^*(U)F_j^{a,\eps}(U)\,d\mu_a(U).
  \label{eq:regulated-gram}
\end{equation}
A normalized quotient block \(q_{a,\eps}\) may be used with Theorem \ref{thm:finite-certificate} only if the following certificate is supplied.

\begin{enumerate}[label=(C\arabic*)]
  \item \textbf{Observable identity.}  The observables \(F_i^{a,\eps}\) are gauge-invariant Wilson observables in the same regulated theory whose reflection positivity is being studied.
  \item \textbf{Reflection identity.}  The reflection \(\Theta\) and positive-time algebra are the ones used in the intended OS reconstruction.
  \item \textbf{Normalization legitimacy.}  The operation producing \(q_{a,\eps}\) from \(N_{a,\eps}\) is implemented by a positivity-preserving same-theory operation: an invertible congruence, a positive factorization, an operator renormalization, or a proved equivalent construction.  Pair-dependent entrywise division is not legitimate without such a proof.
  \item \textbf{Coefficient extraction.}  The expansion
  \begin{equation}
    q_{a,\eps}=\one\one^*+\eps H_a+\eps^2K_a+r_{a,\eps}
    \label{eq:a-eps-expansion}
  \end{equation}
  is proved in operator norm for the stated finite block.
  \item \textbf{Uniform remainders.}  The remainder \(r_{a,\eps}\), endpoint terms, corner terms, and boundary-OPE truncation errors are controlled in the cutoff/collar regime used by the application.
  \item \textbf{No hidden spectral input.}  The estimates do not assume the mass gap or exponential clustering they are intended to help establish.
\end{enumerate}

These conditions are not caveats appended to the finite theorem.  They are the bridge that turns a finite Hermitian-matrix criterion into a legitimate Wilson/OS calculation.  Once they are met, the theorem gives either a positive finite margin for the stated block or a concrete reduced vector showing where the proposed positivity step requires additional work.

\begin{corollary}[Finite Wilson/OS certificate]
\label{cor:same-theory-certificate}
Assume a regulated Wilson construction supplies a normalized quotient block \(q_{a,\eps}\) satisfying the same-theory certificate above and the expansion \eqref{eq:a-eps-expansion}.  If the construction asserts \(q_{a,\eps}\succeq0\) for small \(\eps>0\), then \(PH_aP\succeq0\) on \(\one^\perp\), and the corrected form
\begin{equation}
  P\bigl(K_a-n^{-1}H_aee^*H_a\bigr)P
  \label{eq:corrected-wilson-form}
\end{equation}
  is positive semidefinite on \(\Ker(PH_aP)\).  A negative vector for either reduced form is a finite diagnostic for that proposed block and normalization.
\end{corollary}

\section{Regulated Wilson data model and result-expansion slots}
\label{sec:expansion-slots}

This section implements the result-expansion obligations needed to turn the finite theorem into a controlled Wilson-block application template.  The replacements below separate proved finite facts, same-theory certificates, coefficient-conversion theorems, and regulator-specific obligations that remain conditional.  They are not assumptions used in Theorem \ref{thm:finite-certificate}.

\subsection{Regulator and observable data}

A concrete Wilson application should begin by fixing one of the following regulated settings.
\begin{enumerate}[label=(R\arabic*)]
  \item \textbf{Wilson plaquette action:}
  \begin{equation}
    d\mu_{a,\beta}(U)
    =Z^{-1}_{a,\beta}\exp\!\left(\frac{\beta}{d_R}\sum_p \Ree\Tr_R U_p\right)
    \prod_{\ell}dU_\ell.
    \label{eq:wilson-action}
  \end{equation}
  \item \textbf{Heat-kernel plaquette action:}
  \begin{equation}
    d\mu_{a,t}(U)
    =Z^{-1}_{a,t}\prod_p K_t(U_p)\prod_{\ell}dU_\ell,
    \label{eq:heat-kernel-action}
  \end{equation}
  where \(K_t\) is the heat kernel on \(G\).
\end{enumerate}
The lattice, reflection plane, positive-time edge convention, boundary collar parameter \(\eps\), loop family \(\gamma_i^\eps\), and representation \(R\) must be part of the data, not inferred from notation.

\begin{proposition}[Finite regulated Wilson data sufficient for the reflected block]
\label{prop:regulator-observable-data}
Fix a finite hypercubic lattice \(\Lambda_a\) with compact gauge group \(G\), normalized Haar measure on each oriented edge, a reflection \(\Theta\) across a time-zero hyperplane, and either a Wilson plaquette action \eqref{eq:wilson-action} or a heat-kernel plaquette action \eqref{eq:heat-kernel-action} for which lattice reflection positivity holds for the chosen reflection plane.  Let \(\mathcal A_{+,a}\) be the gauge-invariant algebra generated by Wilson observables whose unoriented edge support lies in the positive half-lattice or in a declared boundary collar assigned to the positive side.  If
\[
  F_i^{a,\eps}=W_{\gamma_i^\eps,R}\in\mathcal A_{+,a},
  \qquad i=1,\ldots,n,
\]
then the finite reflected matrix
\[
  N_{a,\eps}(i,j)
  =
  \int(\Theta F_i^{a,\eps})^*(U)F_j^{a,\eps}(U)\,d\mu_a(U)
\]
is positive semidefinite.
\end{proposition}

\begin{proof}
The statement is the finite-matrix form of lattice reflection positivity.  For any \(c\in\C^n\), the observable \(F=\sum_i c_iF_i^{a,\eps}\) belongs to \(\mathcal A_{+,a}\).  Reflection positivity for the selected finite lattice measure gives
\[
  0\le \int(\Theta F)^*F\,d\mu_a
  =
  \sum_{i,j}\overline{c_i}c_j
  \int(\Theta F_i^{a,\eps})^*F_j^{a,\eps}\,d\mu_a
  =
  c^*N_{a,\eps}c .
\]
Thus \(N_{a,\eps}\succeq0\).  The proposition deliberately keeps the regulator finite: it supplies the finite Gram block used by Theorem \ref{thm:finite-certificate}, not a continuum reconstruction theorem.
\end{proof}

\begin{remark}[What must be named before using the proposition]
The paper must specify \(G\), the representation \(R\), the plaquette action, the reflection plane, the positive-time edge convention, the collar support convention, and the Wilson paths \(\gamma_i^\eps\).  Without these choices the symbols \(q_{a,\eps}\), \(H_a\), and \(K_a\) do not denote checkable same-theory objects.
\end{remark}

\begin{proposition}[Same-theory normalization certificate]
\label{prop:same-theory-normalization}
Let \(N\succeq0\) be a reflected Gram block obtained from a finite regulated Wilson theory.  A normalized block \(q\) may be used as a same-theory reflected block in Theorem \ref{thm:finite-certificate} if one of the following verifiable certificates is supplied:
\begin{enumerate}[label=(Q\arabic*)]
  \item \textbf{One-body congruence.}  There is an invertible diagonal or block-diagonal matrix \(D\), obtained from nonzero one-body normalization factors of the same observables, such that \(q=DND^*\).
  \item \textbf{Positive-kernel factorization.}  There is a positive semidefinite kernel \(L\succeq0\), with \(q=N\Had L\), and \(L\) is realized by an auxiliary same-theory positive factor or by an explicitly stated positive kernel construction.
  \item \textbf{Direct Gram realization.}  There are normalized observables \(\widetilde F_i\in\mathcal A_{+,a}\) in the same regulated theory such that
  \[
    q_{ij}=\int(\Theta\widetilde F_i)^*\widetilde F_j\,d\mu_a .
  \]
\end{enumerate}
In each case \(q\succeq0\).  Pair-dependent entrywise division \(N\odiv B\) is not a certificate unless it is reduced to one of the three cases above.
\end{proposition}

\begin{proof}
In case (Q1), \(x^*q x=(D^*x)^*N(D^*x)\ge0\).  In case (Q2), the Schur product theorem gives \(N\Had L\succeq0\).  In case (Q3), positivity follows directly from reflection positivity applied to \( \sum_i c_i\widetilde F_i\).  These arguments are finite-dimensional and internal to the regulated theory once the stated objects are named.
\end{proof}

\begin{remark}[Why an entrywise quotient is not enough]
The condition is necessary as a proof obligation, not as a stylistic preference.  Appendix \ref{lem:app-entrywise-division} gives positive semidefinite \(N\) and positive-entry \(B\) for which \(N\odiv B\) is indefinite.  Therefore a quotient formula by itself cannot be used as an OS-positive Wilson block.
\end{remark}

\subsection{Coefficient extraction slots}

For a fixed finite family, a complete application must compute or bound the matrices \(H_a\) and \(K_a\) in \eqref{eq:a-eps-expansion}.  If the calculation is performed logarithmically, then one must convert log-coefficients into matrix coefficients before applying Theorem \ref{thm:finite-certificate}.  Appendix \ref{app:quotient-algebra} records the finite conversion.

\begin{proposition}[Coefficient extraction from a normalized finite block]
\label{prop:HK-extraction}
Fix a finite same-theory normalized block \(q_{a,\eps}\) satisfying Proposition \ref{prop:same-theory-normalization}.  Suppose that, uniformly over the finite index set,
\[
  \log(q_{a,\eps})_{ij}
  =
  \eps h_{a,ij}
  +
  \eps^2\ell_{a,ij}
  +
  \eta_{a,ij}(\eps),
  \qquad
  \max_{i,j}|\eta_{a,ij}(\eps)|\le R_a(\eps),
  \qquad
  R_a(\eps)=o(\eps^2).
\]
Assume the branch of the logarithm is chosen in a neighborhood of \(1\), and \(h_a=h_a^*\), \(\ell_a=\ell_a^*\).  Then
\[
  q_{a,\eps}
  =
  \one\one^*
  +\eps H_a
  +\eps^2K_a
  +r_{a,\eps},
  \qquad
  H_a=h_a,
  \qquad
  K_a=\ell_a+\frac12h_a\Had h_a ,
\]
with \(\norm{r_{a,\eps}}\le C_n(R_a(\eps)+O_a(\eps^3))\) for a finite-dimensional norm constant \(C_n\).  In particular, the \(H_a,K_a\) entering Theorem \ref{thm:finite-certificate} are matrix coefficients of \(q_{a,\eps}\), not merely logarithmic OPE coefficients.
\end{proposition}

\begin{proof}
For each pair \((i,j)\),
\[
  \exp\{\eps h_{a,ij}+\eps^2\ell_{a,ij}+\eta_{a,ij}(\eps)\}
  =
  1+\eps h_{a,ij}
  +\eps^2\left(\ell_{a,ij}+\frac12h_{a,ij}^2\right)
  +O_a(\eps^3)+O(R_a(\eps)).
\]
Since the index set is finite, the entrywise bound gives the stated operator-norm bound after changing constants.  Hermiticity follows from the Hermiticity of \(q_{a,\eps}\) and of the coefficient matrices.
\end{proof}

\begin{remark}[Regulator-specific status]
This proposition converts already-derived logarithmic coefficients into the matrix coefficients required by the finite theorem.  It does not derive \(h_a\) or \(\ell_a\) from a particular lattice action.  A completed Wilson application must still compute those quantities from the chosen regulator and observable family.
\end{remark}

\begin{proposition}[Concrete finite three-observable coefficient certificate]
\label{prop:three-observable-certificate}
Let
\[
  e=\frac1{\sqrt3}(1,1,1)^T,\qquad
  u=\frac1{\sqrt2}(1,-1,0)^T,\qquad
  z=\frac1{\sqrt6}(1,1,-2)^T,
\]
and let \(P=I-ee^*\).  Define finite Hermitian coefficient matrices
\[
  H=2uu^*+ez^*+ze^*,
  \qquad
  K=zz^* .
\]
Then \(PHP\) is strictly positive on \(\Span\{u\}\), has \(\Ker(PHP)=\Span\{z\}\) inside \(\one^\perp\), and
\[
  z^*P\bigl(K-3^{-1}Hee^*H\bigr)Pz=\frac23>0.
\]
Consequently the finite family
\[
  q_\eps=\one\one^*+\eps H+\eps^2K
\]
passes the strict Schur--Feshbach validation test for all sufficiently small \(\eps>0\).
\end{proposition}

\begin{proof}
The vectors \(e,u,z\) are orthonormal, \(P\) is the projection onto \(\Span\{u,z\}\), and
\[
  PHP=2uu^*.
\]
Thus the first-order reduced form has margin \(2\) on \(\Span\{u\}\) and kernel \(\Span\{z\}\).  Moreover \(He=z\), so
\[
  z^*P(K-3^{-1}Hee^*H)Pz
  =
  z^*Kz-\frac13|e^*Hz|^2
  =
  1-\frac13
  =
  \frac23.
\]
Proposition \ref{prop:sufficiency} gives positivity for sufficiently small \(\eps\).
\end{proof}

\begin{remark}[Interpretation of the finite outcome]
This is a genuine finite Schur--Feshbach outcome: it has a named block size, a first-order null sector, and a strict corrected second-order scalar.  It is not yet a regulator-specific Wilson calculation.  To promote it to a Wilson result, the matrices \(H,K\) above must be derived as \(H_a,K_a\) from a same-theory three-hairpin block.
\end{remark}

\subsection{Endpoint and remainder slots}

Endpoint/corner terms should not be compressed into unexplained constants.  Color factors, path-ordering factors, representation normalizations, and incidence weights must be derived in the regulator that supplies \(q_{a,\eps}\).

\begin{proposition}[Endpoint prefactor as a same-regulator source derivative]
\label{prop:endpoint-prefactor}
Let \(Z_v(s;b)\) be the finite endpoint source integral in the chosen regulator,
\[
  Z_v(s;b)
  =
  \int_G
  \exp\!\left(
    s\,\sum_\alpha b_\alpha\,\mathcal J_\alpha(g)
  \right)
  d\nu_{v,\mathrm{reg}}(g),
\]
where \(d\nu_{v,\mathrm{reg}}\) is the endpoint conditional group measure induced by the same plaquette action, and \(\mathcal J_\alpha\) are the Lie-algebra source insertions determined by the oriented plaquette sheets and the representation \(R\).  The endpoint fourth-cumulant contribution has the form
\[
  C_{\mathrm{col}}^{(G,R,\mathrm{reg})}
  =
  \frac{1}{4!}
  \sum_{\alpha,\beta,\gamma,\delta}
  \mathcal I_{\alpha\beta\gamma\delta}^{\mathrm{path}}
  \left.
  \partial_{\alpha}\partial_{\beta}\partial_{\gamma}\partial_{\delta}
  \log Z_v(s;b)
  \right|_{s=0},
\]
where \(\mathcal I^{\mathrm{path}}\) is the explicitly counted incidence tensor of the four source insertions at the endpoint.  This formula is the acceptable replacement for an unexplained numerical endpoint constant.
\end{proposition}

\begin{proof}
Expanding the path-ordered Wilson factors at the endpoint gives a finite polynomial in source insertions before integration.  The connected fourth-order contribution is obtained by differentiating \(\log Z_v\) four times at the origin.  The factor \(1/4!\) and the tensor \(\mathcal I^{\mathrm{path}}\) account respectively for the Taylor coefficient and for the ordered incidence multiplicities.  All representation and orientation dependence is therefore contained in \(\mathcal J_\alpha\), \(d\nu_{v,\mathrm{reg}}\), and \(\mathcal I^{\mathrm{path}}\).
\end{proof}

\begin{remark}[Promotion condition]
The paper may use \(C_{\mathrm{col}}^{(G,R,\mathrm{reg})}\) in a claim only after \(d\nu_{v,\mathrm{reg}}\), the source insertions, the trace normalization, and the incidence tensor are written for the specific regulator.  Until then the endpoint scalar remains conditional algebra.
\end{remark}

\begin{proposition}[Uniform endpoint remainder criterion]
\label{prop:uniform-endpoint-remainder}
Fix a finite Wilson block and a compact source domain \(\mathcal S\) for all endpoint source variables used in that block.  Suppose every endpoint log-partition function \(\ell_v(s,z)=\log Z_v(sz)\) is \(C^6\) on \(|s|\le2h\), \(z\in\mathcal S\), and
\[
  M_{6}:=
  \sup_{v}\sup_{z\in\mathcal S}\sup_{|s|\le2h}
  \left|\frac{d^6}{ds^6}\ell_v(s,z)\right|
  <\infty .
\]
Then the five-point fourth-difference certificate satisfies
\[
  \left|D_h\ell_v(z)-\frac{d^4}{ds^4}\ell_v(0,z)\right|
  \le \frac{h^2}{6}M_6
\]
uniformly in \(v\) and \(z\).  If, in addition, the omitted endpoint and collar-tail terms obey
\[
  |E_{\mathrm{tail}}(a,\eps)|+\norm{r_{a,\eps}}\le \rho(a,\eps),
  \qquad
  \rho(a,\eps)=o(\eps^2),
\]
in the same cutoff/collar regime, then the endpoint estimate is small enough to be compared with the finite Schur--Feshbach margin.
\end{proposition}

\begin{proof}
The fourth-difference estimate is the standard centered five-point formula with sixth-derivative remainder, applied uniformly over the finite endpoint set and compact source domain.  The second displayed inequality is exactly the operator-norm hypothesis needed to preserve the coefficient expansion \(q_{a,\eps}=\one\one^*+\eps H_a+\eps^2K_a+o(\eps^2)\).  Since the finite block has fixed dimension, the endpoint and tail bounds enter the same norm used by Theorem \ref{thm:finite-certificate}.
\end{proof}

\begin{remark}[Demotion boundary]
If the source domain, sixth-derivative bound, or tail estimate is not supplied in the same regulator as \(q_{a,\eps}\), endpoint algebra should remain outside the paper's result claim.
\end{remark}

\begin{remark}[Prior-art boundary and novelty claim]
\label{rem:prior-art-boundary}
The finite ingredients used here have recognizable ancestors.  Schur complements and Feshbach maps are standard finite-dimensional reduction tools \cite{Feshbach1958,Feshbach1962}; Osterwalder--Schrader positivity is the reconstruction positivity condition \cite{OsterwalderSchrader1973,OsterwalderSchrader1975}; lattice gauge reflection positivity and transfer-matrix positivity are classical parts of constructive lattice gauge theory \cite{OsterwalderSeiler1978,Luscher1977,Seiler1982}; and Wilson loops are the standard gauge-invariant observables introduced by Wilson \cite{Wilson1974}.

The novelty claimed by this paper is therefore bounded.  It is not the Schur complement itself, not reflection positivity itself, and not a construction of Yang--Mills theory.  The contribution is the packaging of the rank-one reflected-block collision into the explicit finite coefficient test
\[
  PHP,\qquad
  P(K-n^{-1}Hee^*H)P,
\]
together with a same-theory quotient warning and a reproducible finite model showing how a first-order null sector is resolved by the corrected second-order form.  Any stronger novelty claim requires a regulator-specific Wilson coefficient calculation that discharges Propositions \ref{prop:regulator-observable-data}--\ref{prop:uniform-endpoint-remainder}.
\end{remark}

\section{Three-hairpin block as an application template}
\label{sec:three-hairpin-template}

A three-observable block is the smallest common setting where the second-order null-sector test can be one-dimensional.  Suppose \(n=3\) and the first reduced form \(A=PHP\) has a one-dimensional null sector represented by
\begin{equation}
  z=(1,1,-2)^T,
  \qquad z\perp\one.
  \label{eq:z-vector}
\end{equation}
Then the second finite criterion reduces to
\begin{equation}
  z^*Cz\ge0,
  \qquad
  C=P\bigl(K-3^{-1}Hee^*H\bigr)P.
  \label{eq:three-hairpin-criterion}
\end{equation}
If the scalar is positive, define the finite margin
\begin{equation}
  M_G:=z^*Cz>0.
  \label{eq:MG}
\end{equation}
If it is negative, the block supplies a concrete finite diagnostic.  In either case, the calculation is meaningful only after the same-theory transfer certificate in Section \ref{sec:same-theory} is met.

A regulator-specific three-hairpin Wilson section must eventually contain the following data.  The present implementation supplies an exact finite coefficient model and marks the remaining Wilson-path derivation boundary.
\begin{enumerate}[label=(T\arabic*)]
  \item the three Wilson paths \(\gamma_1^\eps,\gamma_2^\eps,\gamma_3^\eps\), drawn or specified combinatorially;
  \item their common reflection-plane footprint and the local geometric differences between them;
  \item the reflected Gram block \(N_{a,\eps}\);
  \item the legitimate same-theory normalization producing \(q_{a,\eps}\);
  \item explicit or bounded coefficients \(H_a,K_a\);
  \item the values of \(PH_aP\) on a basis of \(\one^\perp\);
  \item the scalar \(z^*P(K_a-3^{-1}H_aee^*H_a)Pz\), if \(z\in\Ker(PH_aP)\);
  \item uniform collar/cutoff remainders.
\end{enumerate}

\begin{center}
\begin{tabular}{p{0.16\linewidth}p{0.20\linewidth}p{0.20\linewidth}p{0.22\linewidth}}
\toprule
Pair/channel & First-order contribution & Corrected second-order role & Status \\
\midrule
\((u,u)\) & \(u^*PHPu=2\) & not needed for the null test & strict first-order margin \\
\((z,z)\) & \(z^*PHPz=0\) & \(z^*P(K-3^{-1}Hee^*H)Pz=2/3\) & strict second-order null-sector margin \\
\((u,z)\) & \(u^*PHPz=0\) & \(u^*P(K-3^{-1}Hee^*H)Pz=0\) & no mixed obstruction in this model \\
\bottomrule
\end{tabular}
\end{center}

In the standard coordinate basis this coefficient model is
\[
H=
2uu^*+ez^*+ze^*,
\qquad
K=zz^*,
\]
with \(e,u,z\) as in Proposition \ref{prop:three-observable-certificate}.  The table is reproducible from the exact certificate generated in
\[
\texttt{three\_hairpin\_certificate.json}
\]
in the result-elaboration report directory.
As with Proposition \ref{prop:three-observable-certificate}, this table is a finite coefficient realization of the three-hairpin algebra.  It becomes a Wilson three-hairpin table only after the same entries are derived from named Wilson paths in a named regulator.

\section{Endpoint cumulant extension: conditional algebra and placeholders}
\label{sec:endpoint-extension}

This section records a conditional algebraic template for a post-\(H,K\) endpoint estimate.  It is not used in the proof of the finite theorem and should not be read as a completed Wilson calculation.

Assume a three-hairpin calculation has produced a positive finite margin \(M_G\).  Let \(R\) denote the next connected boundary-polarization contribution.  In the null channel \(z=(1,1,-2)^T\), the pair weights are
\begin{equation}
  z^*Rz=2R_{12}-4R_{13}-4R_{23}.
  \label{eq:null-weights}
\end{equation}
For endpoint boxes \(v\), suppose a same-regulator calculation gives connected fourth source derivatives
\begin{equation}
  \kappa_{4,v}^{\mathrm{sing}}(b)
  =\left.\partial_s^4\log Z_v(s;b)\right|_{s=0}.
  \label{eq:kappa4}
\end{equation}
Here \(Z_v(s;b)\) must be defined as an actual finite group integral or conditional endpoint partition function in the chosen regulator.

With a derived prefactor \(C_{\mathrm{col}}^{(G,R,\mathrm{reg})}\), the endpoint scalar has the conditional form
\begin{equation}
  R_{\mathrm{null}}^{\mathrm{end}}
  =C_{\mathrm{col}}^{(G,R,\mathrm{reg})}
  \sum_v
  \{2\kappa_{4,v}(b_{12,v})-4\kappa_{4,v}(b_{13,v})-4\kappa_{4,v}(b_{23,v})\}
  +E_{\mathrm{tail}}.
  \label{eq:endpoint-scalar}
\end{equation}
If
\begin{equation}
  \kappa_{4,v}(b_{ij,v})=-\chi_v B_{ij,v}+\epsilon_{ij,v},
  \qquad \chi_v\ge0,
  \label{eq:kappa-expansion}
\end{equation}
and
\begin{equation}
  D_v=2B_{12,v}-4B_{13,v}-4B_{23,v},
  \label{eq:Dv}
\end{equation}
then the purely algebraic lower bound is
\begin{equation}
  \snull
  \ge
  M_G
  -C_{\mathrm{col}}^{(G,R,\mathrm{reg})}\sum_v\chi_v[D_v]_+
  -C_{\mathrm{col}}^{(G,R,\mathrm{reg})}\sum_{v,ij}|w_{ij}|\,|\epsilon_{ij,v}|
  -|E_{\mathrm{tail}}|,
  \label{eq:safe-endpoint-bound}
\end{equation}
where \((w_{12},w_{13},w_{23})=(2,-4,-4)\).  Appendix \ref{app:endpoint-certificates} proves this algebraic implication.  The remaining regulator-specific QFT work is exactly the content bounded by Propositions \ref{prop:endpoint-prefactor} and \ref{prop:uniform-endpoint-remainder}: define \(Z_v\), derive the prefactor, compute incidence data, and bound the tails in the same regulator.

A source-derivative computation may use the five-point fourth-difference stencil
\begin{align}
D_h\ell_v(z)
&=\frac{\ell_v(2hz)-4\ell_v(hz)+6\ell_v(0)-4\ell_v(-hz)+\ell_v(-2hz)}{h^4},\notag\\
\left|D_h\ell_v(z)-K_v^{\mathrm{sing}}[z]\right|
&\le\frac{h^2}{6}M_{6,v}(z,h),
\label{eq:central-difference}
\end{align}
where \(\ell_v(s,z)=\log Z_v(sz)\).  This is a numerical or symbolic certificate only after \(M_{6,v}\) is uniformly bounded in the stated source domain.

\section{Discussion: what has been established and what remains}

The paper establishes a finite Schur--Feshbach tangent-cone criterion.  Its central formula is
\[
  F_\eps=
  \eps PHP+
  \eps^2P\bigl(K-n^{-1}Hee^*H\bigr)P+o(\eps^2).
\]
This formula is independent of Yang--Mills theory.  It is relevant to Wilson/OS constructions because normalized reflected Wilson blocks can have rank-one collision limits, and because the finite coefficient test is unavoidable once such a limit is present.

The finite theorem advances the programme in four ways.
\begin{enumerate}[label=(A\arabic*)]
  \item It turns a singular reflection-positivity step into a finite coefficient calculation.
  \item It identifies the corrected second-order form, preventing the common mistake of testing only \(PKP\).
  \item It separates a positive finite margin from a finite diagnostic vector, so a Wilson calculation can either support a construction step or localize the next required repair.
  \item It states the same-theory transfer certificate needed before a quotient block can be interpreted as a Wilson/OS object.
\end{enumerate}

The paper does not construct four-dimensional Yang--Mills theory, does not prove OS reconstruction for a continuum gauge field, and does not prove a mass gap.  The result-bearing expansion implemented here replaces the slots by finite certificates, coefficient-conversion theorems, an exact three-observable Schur--Feshbach model, and explicit demotion boundaries.  A future regulator-specific Wilson calculation must still derive the same coefficient data from named Wilson paths.

\appendix

\section{Conceptual explainer for QFT and geometry readers}
\label{app:conceptual-explainer}

\subsection{Reflection positivity in finite form}

In Euclidean QFT, one constructs imaginary-time correlation functions and then reconstructs a Hilbert-space theory in real time.  Reflection positivity is the condition that makes the inner product nonnegative after reflecting positive-time observables.  For finitely many observables \(F_i\), it is simply the statement that the matrix \(N_{ij}=S((\Theta F_i)^*F_j)\) is positive semidefinite.

\subsection{Why Wilson loops enter}

In gauge theory, the gauge potential itself is not gauge invariant.  Wilson loops, built from traces of holonomies, are gauge invariant.  They are therefore natural probes for a gauge-invariant Euclidean construction.  When Wilson loops approach a reflection plane, their reflected overlaps can become almost identical at leading order.  This is the finite rank-one limit \(q_\eps\to\one\one^*\).

\subsection{Why the Schur complement is forced}

The vector \(\one\) represents the common leading mode.  Quotient vectors in \(\one^\perp\) can still couple into this large rank-one direction at first order and return at second order.  Removing \(\one\) therefore requires a Schur complement, not ordinary projection.  That is why the second-order coefficient contains
\[
  -n^{-1}PHee^*HP.
\]
This term is the finite trace of coupling through the rank-one direction.

\subsection{What the theorem contributes}

Given \(q_\eps=\one\one^*+\eps H+\eps^2K+o(\eps^2)\), the theorem says to compute
\[
  A=PHP,
  \qquad
  C=P(K-n^{-1}Hee^*H)P.
\]
If \(A\) is positive with margin, the finite block has a first-order positivity margin.  If \(A\) has a kernel, test \(C\) on that kernel.  A negative vector in either test locates a finite defect in the proposed positivity step for that block.  A positive reduced margin gives a quantity against which endpoint and remainder estimates can be compared.

\subsection{Relation to Yang--Mills}

The Yang--Mills mass-gap problem asks for much more: construction of a nontrivial four-dimensional quantum Yang--Mills theory and proof of a positive spectral gap \cite{JaffeWitten2000,ClayYM}.  The finite theorem here is an earlier local ingredient.  It tells a Wilson/OS reconstruction programme how to audit finite normalized reflected blocks before using them in a continuum argument.

\section{Expanded proofs of the finite criterion}
\label{app:expanded-proofs}

This appendix expands the finite-dimensional arguments.  No gauge-theoretic input is used.

\subsection{Scalar Schur complement}

\begin{lemma}[Scalar Schur complement]
\label{lem:app-scalar-schur}
Let \(E\) be a finite-dimensional complex Hilbert space, let \(a>0\), let \(b\in E\), and let \(D=D^*\in\operatorname{End}(E)\).  On \(\C\oplus E\), define
\[
  M=\begin{pmatrix}a&b^*\\ b&D\end{pmatrix}.
\]
Then
\[
  M\succeq0\quad\Longleftrightarrow\quad D-a^{-1}bb^*\succeq0.
\]
\end{lemma}

\begin{proof}
For \((s,y)\in\C\oplus E\), completing the square gives
\begin{align*}
  \begin{pmatrix}\bar s&y^*\end{pmatrix}
  M
  \begin{pmatrix}s\\y\end{pmatrix}
  &=a|s|^2+2\Ree(\bar s b^*y)+y^*Dy\\
  &=a\left|s+a^{-1}b^*y\right|^2+y^*(D-a^{-1}bb^*)y.
\end{align*}
The equivalence follows by first assuming \(D-a^{-1}bb^*\succeq0\), and conversely by choosing \(s=-a^{-1}b^*y\).
\end{proof}

\subsection{Expansion of the reduced block}

With respect to \(\C^n=\C e\oplus P\C^n\), write
\[
  q_\eps=\begin{pmatrix}
    \lambda_\eps & b_\eps^*\\
    b_\eps & D_\eps
  \end{pmatrix},
  \qquad
  \lambda_\eps=e^*q_\eps e,
  \quad b_\eps=Pq_\eps e,
  \quad D_\eps=Pq_\eps P.
\]
The leading term \(nee^*\) contributes only to the scalar block.  Hence
\begin{align*}
  \lambda_\eps
  &=n+\eps e^*He+\eps^2e^*Ke+o(\eps^2),\\
  b_\eps
  &=\eps PHe+\eps^2PKe+o(\eps^2),\\
  D_\eps
  &=\eps PHP+\eps^2PKP+o(\eps^2).
\end{align*}
Thus \(\lambda_\eps=n+O(\eps)>0\), and \(\lambda_\eps^{-1}=n^{-1}+O(\eps)\).

\begin{lemma}[Vacuum Schur--Feshbach expansion]
\label{lem:app-feshbach-expansion}
The Schur complement
\[
  F_\eps=Pq_\eps P-Pq_\eps e\,(e^*q_\eps e)^{-1}\,e^*q_\eps P
\]
satisfies
\[
  F_\eps=\eps PHP+
  \eps^2P\bigl(K-n^{-1}Hee^*H\bigr)P+o(\eps^2).
\]
\end{lemma}

\begin{proof}
Using the block expansions,
\begin{align*}
  Pq_\eps e\,(e^*q_\eps e)^{-1}\,e^*q_\eps P
  &=\bigl(\eps PHe+O(\eps^2)\bigr)
    \bigl(n^{-1}+O(\eps)\bigr)
    \bigl(\eps e^*HP+O(\eps^2)\bigr)\\
  &=\eps^2n^{-1}PHee^*HP+o(\eps^2).
\end{align*}
Subtracting from \(Pq_\eps P=\eps PHP+\eps^2PKP+o(\eps^2)\) gives the claim.
\end{proof}

\subsection{PSD tangent-cone necessity}

\begin{lemma}[Two-order PSD perturbation test]
\label{lem:app-two-order-test}
Let \(E\) be finite-dimensional and let
\[
  G_\eps=\eps A+\eps^2C+r_\eps,
  \qquad \norm{r_\eps}=o(\eps^2),
\]
where \(A=A^*\), \(C=C^*\).  If \(G_\eps\succeq0\) for all sufficiently small \(\eps>0\), then \(A\succeq0\).  If \(\Pi_A\) projects onto \(\Ker A\), then \(\Pi_A C\Pi_A\succeq0\) on \(\Ker A\).
\end{lemma}

\begin{proof}
For fixed \(v\in E\),
\[
  0\le v^*G_\eps v
  =\eps v^*Av+\eps^2v^*Cv+o(\eps^2).
\]
Division by \(\eps\) gives \(v^*Av\ge0\).  If \(v\in\Ker A\), division by \(\eps^2\) gives \(v^*Cv\ge0\).
\end{proof}

\subsection{Strict sufficiency}

\begin{lemma}[Strict two-block validation]
\label{lem:app-strict-validation}
Let \(F_\eps=\eps A+\eps^2C+o(\eps^2)\) on a finite-dimensional Hilbert space \(E\), with \(A\succeq0\).  Let \(\Pi\) project onto \(\Ker A\) and let \(Q=I_E-\Pi\).  If
\[
  QAQ\succeq\alpha Q,
  \qquad
  \Pi C\Pi\succeq\beta\Pi
\]
for some \(\alpha,\beta>0\), then \(F_\eps\succeq0\) for all sufficiently small \(\eps>0\).
\end{lemma}

\begin{proof}
The \(Q\)-block is \(QF_\eps Q\succeq(\alpha/2)\eps Q\) for small \(\eps\).  Since \(A\Pi=0\), the off-diagonal block is \(\Pi F_\eps Q=O(\eps^2)\).  The Schur complement of the \(Q\)-block is
\[
  \Pi F_\eps\Pi-
  \Pi F_\eps Q(QF_\eps Q)^{-1}QF_\eps\Pi
  =\eps^2\Pi C\Pi+o(\eps^2)-O(\eps^3),
\]
which is positive for small \(\eps\).  The Schur complement criterion gives \(F_\eps\succeq0\).
\end{proof}

\section{Auxiliary algebra for quotient blocks}
\label{app:quotient-algebra}

\begin{lemma}[One-body factors are congruences]
\label{lem:app-one-body}
Let \(r=r^*\), and let \(D=\diag(d_1,\ldots,d_n)\) be invertible.  If \(q=DrD^*\), then \(q\succeq0\) if and only if \(r\succeq0\).
\end{lemma}

\begin{proof}
For every \(x\), \(x^*qx=(D^*x)^*r(D^*x)\).  Since \(D^*\) is invertible, the equivalence follows.
\end{proof}

\begin{lemma}[Zero-sum cancellation of additive one-body logarithms]
\label{lem:app-zero-sum}
For \(a\in\C^n\), set \(L=a\one^*+\one a^*\).  Then \(v^*Lv=0\) for all \(v\in\one^\perp\).
\end{lemma}

\begin{proof}
If \(v\in\one^\perp\), then \(\one^*v=0\) and \(v^*\one=0\).  Thus
\[
  v^*Lv=(v^*a)(\one^*v)+(v^*\one)(a^*v)=0.
\]
\end{proof}

\begin{lemma}[From log coefficients to matrix coefficients]
\label{lem:app-log-HK}
Assume, entrywise and uniformly over a finite index set, that
\[
  \log (q_\eps)_{ij}=\eps h_{ij}+\eps^2\ell_{ij}+o(\eps^2),
\]
with \(h=h^*\), \(\ell=\ell^*\).  Then
\[
  q_\eps=\one\one^*+\eps H+\eps^2K+o(\eps^2),
  \qquad
  H=h,
  \qquad
  K=\ell+\frac12 h\Had h,
\]
where \(\Had\) denotes entrywise product.
\end{lemma}

\begin{proof}
For each entry, \(e^{\eps h_{ij}+\eps^2\ell_{ij}+o(\eps^2)}=1+\eps h_{ij}+\eps^2(\ell_{ij}+h_{ij}^2/2)+o(\eps^2)\).  Since the index set is finite, entrywise uniformity gives the matrix-norm expansion.
\end{proof}

\begin{lemma}[Entrywise division can destroy positivity]
\label{lem:app-entrywise-division}
There exist positive semidefinite matrices \(N\) and \(B\), with every entry of \(B\) strictly positive, such that \(N\odiv B\) is indefinite.
\end{lemma}

\begin{proof}
Take
\[
  N=\begin{pmatrix}1&1\\1&1\end{pmatrix},
  \qquad
  B=\begin{pmatrix}1&1/10\\1/10&1\end{pmatrix}.
\]
The eigenvalues of \(N\) are \(2,0\), and the eigenvalues of \(B\) are \(11/10,9/10\).  But
\[
  N\odiv B=\begin{pmatrix}1&10\\10&1\end{pmatrix}
\]
has eigenvalues \(11\) and \(-9\).
\end{proof}

\section{Endpoint algebra}
\label{app:endpoint-certificates}

\begin{lemma}[Endpoint lower bound]
\label{lem:app-endpoint-bound}
Assume
\[
  R_{\mathrm{null}}^{\mathrm{end}}
  =C\sum_v\{2\kappa_{4,v}(b_{12,v})-4\kappa_{4,v}(b_{13,v})-4\kappa_{4,v}(b_{23,v})\}
  +E_{\mathrm{tail}},
\]
where \(C\ge0\).  Suppose
\[
  \kappa_{4,v}(b_{ij,v})=-\chi_vB_{ij,v}+\epsilon_{ij,v},
  \qquad \chi_v\ge0,
\]
and define
\[
  D_v=2B_{12,v}-4B_{13,v}-4B_{23,v},
  \qquad (w_{12},w_{13},w_{23})=(2,-4,-4).
\]
Then
\[
  M_G+R_{\mathrm{null}}^{\mathrm{end}}
  \ge
  M_G-C\sum_v\chi_v[D_v]_+-C\sum_{v,ij}|w_{ij}|\,|\epsilon_{ij,v}|-|E_{\mathrm{tail}}|.
\]
\end{lemma}

\begin{proof}
Substitution gives
\[
  R_{\mathrm{null}}^{\mathrm{end}}
  =C\sum_v\left\{-\chi_vD_v+\sum_{ij}w_{ij}\epsilon_{ij,v}\right\}+E_{\mathrm{tail}}.
\]
Since \(\chi_v\ge0\), \(-\chi_vD_v\ge-\chi_v[D_v]_+\).  Also,
\[
  \sum_{ij}w_{ij}\epsilon_{ij,v}\ge-\sum_{ij}|w_{ij}|\,|\epsilon_{ij,v}|,
  \qquad E_{\mathrm{tail}}\ge-|E_{\mathrm{tail}}|.
\]
The result follows by summing and adding \(M_G\).
\end{proof}

\begin{lemma}[Fourth derivative from the five-point stencil]
\label{lem:app-fourth-difference}
Let \(f\in C^6([-2h,2h])\), and set
\[
  D_hf=
  \frac{f(2h)-4f(h)+6f(0)-4f(-h)+f(-2h)}{h^4}.
\]
Then
\[
  \left|D_hf-f^{(4)}(0)\right|
  \le
  \frac{h^2}{6}\sup_{|s|\le2h}|f^{(6)}(s)|.
\]
\end{lemma}

\begin{proof}
Taylor expand \(f(\pm h)\) and \(f(\pm2h)\) about zero with integral or Lagrange remainder.  The coefficients of \(f,f',f'',f''',f^{(5)}\) cancel, the coefficient of \(f^{(4)}(0)\) is one after division by \(h^4\), and the sixth-order remainder is bounded by the displayed constant.  Equivalently, this is the standard centered five-point fourth-derivative formula.
\end{proof}

\section{Author-facing result-expansion checklist}
\label{app:author-checklist}

This appendix is intended for draft development.  It distinguishes expressions that belong in a final paper from placeholders that must be expanded or removed before submission.

\begin{center}
\begin{tabular}{p{0.30\linewidth}p{0.58\linewidth}}
\toprule
Keep in final theorem & \(q_\eps=\one\one^*+\eps H+\eps^2K+o(\eps^2)\), \(e=n^{-1/2}\one\), \(P=I-ee^*\), the Schur complement \(F_\eps\), the coefficients \(PHP\) and \(P(K-n^{-1}Hee^*H)P\), and the strict-margin proposition. \\
\midrule
Keep in Wilson application & The specified regulator, reflection plane, Wilson observables, Gram block, same-theory normalization, coefficient extraction, finite sign test, and uniform remainder estimates. \\
\midrule
Keep only conditionally & Endpoint cumulants, incidence weights, and color/combinatorial prefactors, unless each is derived in the same regulator. \\
\midrule
Remove before journal submission & Any operational provenance, scoring traces, or claims that the finite theorem alone proves continuum Yang--Mills reflection positivity or a mass gap. \\
\bottomrule
\end{tabular}
\end{center}

\begin{thebibliography}{99}

\bibitem{YangMills1954}
C.~N. Yang and R.~L. Mills.
\newblock Conservation of isotopic spin and isotopic gauge invariance.
\newblock \emph{Physical Review} 96, 191--195, 1954.

\bibitem{OsterwalderSchrader1973}
K.~Osterwalder and R.~Schrader.
\newblock Axioms for Euclidean Green's functions.
\newblock \emph{Communications in Mathematical Physics} 31, 83--112, 1973.
\newblock doi:10.1007/BF01645738.

\bibitem{OsterwalderSchrader1975}
K.~Osterwalder and R.~Schrader.
\newblock Axioms for Euclidean Green's functions II.
\newblock \emph{Communications in Mathematical Physics} 42, 281--305, 1975.
\newblock doi:10.1007/BF01608978.

\bibitem{Wilson1974}
K.~G. Wilson.
\newblock Confinement of quarks.
\newblock \emph{Physical Review D} 10, 2445--2459, 1974.
\newblock doi:10.1103/PhysRevD.10.2445.

\bibitem{OsterwalderSeiler1978}
K.~Osterwalder and E.~Seiler.
\newblock Gauge field theories on a lattice.
\newblock \emph{Annals of Physics} 110, 440--471, 1978.

\bibitem{Luscher1977}
M.~L\"uscher.
\newblock Construction of a selfadjoint, strictly positive transfer matrix for Euclidean lattice gauge theories.
\newblock \emph{Communications in Mathematical Physics} 54, 283--292, 1977.
\newblock doi:10.1007/BF01614090.

\bibitem{Seiler1982}
E.~Seiler.
\newblock \emph{Gauge Theories as a Problem of Constructive Quantum Field Theory and Statistical Mechanics}.
\newblock Springer, 1982.

\bibitem{Feshbach1958}
H.~Feshbach.
\newblock Unified theory of nuclear reactions.
\newblock \emph{Annals of Physics} 5, 357--390, 1958.

\bibitem{Feshbach1962}
H.~Feshbach.
\newblock A unified theory of nuclear reactions II.
\newblock \emph{Annals of Physics} 19, 287--313, 1962.

\bibitem{MenottiOnofri1981}
P.~Menotti and E.~Onofri.
\newblock The action of \(SU(N)\) lattice gauge theory in terms of the heat kernel on the group manifold.
\newblock \emph{Nuclear Physics B} 190, 288--300, 1981.

\bibitem{Polyakov1980}
A.~M. Polyakov.
\newblock Gauge fields as rings of glue.
\newblock \emph{Nuclear Physics B} 164, 171--188, 1980.

\bibitem{DotsenkoVergeles1980}
V.~S. Dotsenko and S.~N. Vergeles.
\newblock Renormalizability of phase factors in the non-Abelian gauge theory.
\newblock \emph{Nuclear Physics B} 169, 527--546, 1980.

\bibitem{BrandtNeriSato1981}
R.~A. Brandt, F.~Neri, and M.~Sato.
\newblock Renormalization of loop functions for all loops.
\newblock \emph{Physical Review D} 24, 879--902, 1981.

\bibitem{KorchemskyRadyushkin1987}
G.~P. Korchemsky and A.~V. Radyushkin.
\newblock Renormalization of the Wilson loops beyond the leading order.
\newblock \emph{Nuclear Physics B} 283, 342--364, 1987.

\bibitem{JaffeWitten2000}
A.~Jaffe and E.~Witten.
\newblock Quantum Yang--Mills theory.
\newblock Clay Mathematics Institute Millennium Problem statement, 2000.
\newblock \url{https://www.claymath.org/wp-content/uploads/2022/06/yangmills.pdf}

\bibitem{ClayYM}
Clay Mathematics Institute.
\newblock Yang--Mills and the Mass Gap.
\newblock \url{https://www.claymath.org/millennium/yang-mills-the-maths-gap/}

\bibitem{Gallier2019}
J.~Gallier.
\newblock The Schur complement and symmetric positive semidefinite matrices.
\newblock Technical note, University of Pennsylvania, 2019.
\newblock \url{https://www.cis.upenn.edu/~jean/schur-comp.pdf}

\end{thebibliography}

\end{document}
